
\documentclass[11pt,a4paper]{article}
\usepackage[T1]{fontenc}
\usepackage{lmodern}
\usepackage{amsmath,amssymb,amsfonts,mathtools}
\usepackage{graphicx,booktabs,longtable,array,tabularx,adjustbox}
\usepackage{geometry,microtype,fancyhdr,lastpage,etoolbox,enumitem}
\usepackage{newunicodechar,textcomp}
\usepackage{hyperref}
\hypersetup{hidelinks,pdfcreator={XeLaTeX / QM1 academic edition}}
\geometry{left=20mm,right=20mm,top=22mm,bottom=23mm}
\setcounter{secnumdepth}{2}
\setlength{\parindent}{0pt}
\setlength{\parskip}{4pt plus 1pt}
\setlength{\emergencystretch}{2em}
\setlength{\tabcolsep}{3pt}
\renewcommand{\arraystretch}{1.16}
\setlist{nosep,topsep=3pt,itemsep=2pt,leftmargin=7mm}
\AtBeginEnvironment{longtable}{\footnotesize\setlength{\tabcolsep}{2.2pt}}
\newcommand{\tightlist}{\setlength{\itemsep}{0pt}\setlength{\parskip}{0pt}}
\providecommand{\Cov}{\operatorname{Cov}}
\providecommand{\Var}{\operatorname{Var}}
\providecommand{\Corr}{\operatorname{Corr}}
\providecommand{\MAD}{\operatorname{MAD}}
\providecommand{\E}{\mathbb{E}}
\providecommand{\R}{\mathbb{R}}
\newunicodechar{—}{\textemdash}
\newunicodechar{–}{\textendash}
\newunicodechar{€}{\texteuro}
\newunicodechar{≤}{\ensuremath{\leq}}
\newunicodechar{≥}{\ensuremath{\geq}}
\newunicodechar{×}{\ensuremath{\times}}
\newunicodechar{÷}{\ensuremath{\div}}
\newunicodechar{Σ}{\ensuremath{\Sigma}}
\newunicodechar{σ}{\ensuremath{\sigma}}
\newunicodechar{μ}{\ensuremath{\mu}}
\newunicodechar{Δ}{\ensuremath{\Delta}}
\newunicodechar{∑}{\ensuremath{\sum}}
\newunicodechar{≈}{\ensuremath{\approx}}
\newunicodechar{√}{\ensuremath{\surd}}
\newunicodechar{≠}{\ensuremath{\ne}}
\newunicodechar{→}{\ensuremath{\to}}
\newunicodechar{∞}{\ensuremath{\infty}}
\newunicodechar{∈}{\ensuremath{\in}}
\newunicodechar{−}{\ensuremath{-}}
\newunicodechar{²}{\ensuremath{^{2}}}
\newunicodechar{³}{\ensuremath{^{3}}}
\newunicodechar{₀}{\ensuremath{_{0}}}
\newunicodechar{₁}{\ensuremath{_{1}}}
\newunicodechar{₂}{\ensuremath{_{2}}}
\newunicodechar{₃}{\ensuremath{_{3}}}
\newunicodechar{₄}{\ensuremath{_{4}}}
\newunicodechar{ₜ}{\ensuremath{_{t}}}
\newunicodechar{ₚ}{\ensuremath{_{p}}}
\newunicodechar{ₙ}{\ensuremath{_{n}}}
\newunicodechar{ₓ}{\ensuremath{_{x}}}
\newunicodechar{ᵢ}{\ensuremath{_{i}}}
\newunicodechar{ŷ}{\ensuremath{\hat y}}
\newunicodechar{Ŷ}{\ensuremath{\widehat Y}}
\newunicodechar{ₘ}{\ensuremath{_{m}}}
\newunicodechar{ₑ}{\ensuremath{_{e}}}
\newunicodechar{ₐ}{\ensuremath{_{a}}}
\newunicodechar{ₖ}{\ensuremath{_{k}}}
\newunicodechar{ₛ}{\ensuremath{_{s}}}
\newunicodechar{ₗ}{\ensuremath{_{l}}}
\newunicodechar{ₕ}{\ensuremath{_{h}}}
\newunicodechar{ₒ}{\ensuremath{_{o}}}
\newunicodechar{₊}{\ensuremath{_{+}}}
\newunicodechar{₋}{\ensuremath{_{-}}}
\newunicodechar{₌}{\ensuremath{_{=}}}
\newunicodechar{₍}{\ensuremath{_{(}}}
\newunicodechar{₎}{\ensuremath{_{)}}}
\newunicodechar{⁻}{\ensuremath{^{-}}}
\newunicodechar{⁺}{\ensuremath{^{+}}}
\newunicodechar{⁰}{\ensuremath{^{0}}}
\newunicodechar{¹}{\ensuremath{^{1}}}
\newunicodechar{⁴}{\ensuremath{^{4}}}
\newunicodechar{⁵}{\ensuremath{^{5}}}
\newunicodechar{⁶}{\ensuremath{^{6}}}
\newunicodechar{⁷}{\ensuremath{^{7}}}
\newunicodechar{⁸}{\ensuremath{^{8}}}
\newunicodechar{⁹}{\ensuremath{^{9}}}
\makeatletter
\renewcommand{\section}{\@startsection{section}{1}{0pt}{2.0ex plus .7ex minus .3ex}{1.0ex}{\large\bfseries}}
\renewcommand{\subsection}{\@startsection{subsection}{2}{0pt}{1.8ex plus .5ex minus .2ex}{.7ex}{\normalsize\bfseries}}
\renewcommand{\subsubsection}{\@startsection{subsubsection}{3}{0pt}{1.3ex plus .4ex minus .2ex}{.5ex}{\normalsize\bfseries}}
\makeatother
\pagestyle{fancy}
\fancyhf{}
\renewcommand{\headrulewidth}{0.35pt}
\renewcommand{\footrulewidth}{0pt}
\fancyhead[L]{\small\textsc{KEDGE Business School --- QM1}}
\fancyhead[R]{\small\textsc{2026--2027}}
\fancyfoot[L]{\small Badr ABOUFARASSE, PhD}
\fancyfoot[R]{\small\thepage\ /\ \pageref{LastPage}}
\setlength{\headheight}{14pt}
\setkeys{Gin}{keepaspectratio}
\begin{document}

\begin{center}
{\small\textsc{KEDGE Business School \quad EBP1 \quad Quantitative Methods 1}}\\[6pt]
{\Large\bfseries Fiche 7 — Cas intégrateurs}\\[6pt]
{\normalsize\bfseries Corrigés intégralement détaillés}\\[6pt]
{\small Badr ABOUFARASSE, PhD \quad --- \quad 2026--2027}
\end{center}
\vspace{2mm}\hrule\vspace{6mm}

\hypertarget{pruxe9paration-aux-uxe9valuations}{%
\section{Préparation aux évaluations}\label{pruxe9paration-aux-uxe9valuations}}

Résoudre des exercices de contrôle continu et de partiel en autonomie.

\emph{Préparation CC / partiel --- 70 min}

\hypertarget{cas-a-pruxe9paration-au-contruxf4le-continu}{%
\subsection{Cas A -- Préparation au contrôle continu}\label{cas-a-pruxe9paration-au-contruxf4le-continu}}

Une école interroge 200 étudiants sur le prix maximal qu'ils accepteraient de payer pour un service optionnel :

\begin{longtable}[]{@{}lllllll@{}}
\toprule\noalign{}
Prix & {\([5,10[\)} & {\([10,15[\)} & {\([15,20[\)} & {\([20,30[\)} & {\([30,40[\)} & {\([40,60[\)} \\
\midrule\noalign{}
\endhead
\bottomrule\noalign{}
\endlastfoot
{\(n_i\)} & 10 & 25 & 45 & 60 & 40 & 20 \\
\end{longtable}

\begin{enumerate}
\item
  Identifier population observée, unité statistique, caractère et nature du caractère.
\item
  Construire le tableau complet : amplitudes, densités, fréquences, ECC, ECD et centres.
\item
  Tracer l'histogramme et les courbes cumulatives.
\item
  Déterminer la classe modale, la médiane, {\(Q_1\)}, {\(Q_3\)} et la moyenne approchée.
\item
  Calculer la variance de population, l'écart-type et le coefficient de variation.
\item
  Déterminer le prix maximal tel qu'environ 65\% des étudiants puissent payer au moins ce prix.
\item
  On demande aussi « satisfait / non satisfait ». Préciser la nature de cette variable et le graphique adapté.
\end{enumerate}

\hypertarget{corriguxe9-duxe9tailluxe9}{%
\subsubsection{Corrigé détaillé}\label{corriguxe9-duxe9tailluxe9}}

Population observée : 200 étudiants interrogés. Unité : un étudiant. Caractère : prix maximal accepté, quantitatif continu.

\begin{longtable}[]{@{}lllllll@{}}
\toprule\noalign{}
Classe & {\([5,10[\)} & {\([10,15[\)} & {\([15,20[\)} & {\([20,30[\)} & {\([30,40[\)} & {\([40,60[\)} \\
\midrule\noalign{}
\endhead
\bottomrule\noalign{}
\endlastfoot
{\(n_i\)} & 10 & 25 & 45 & 60 & 40 & 20 \\
{\(a_i\)} & 5 & 5 & 5 & 10 & 10 & 20 \\
{\(d_i\)} & 2 & 5 & 9 & 6 & 4 & 1 \\
ECC & 10 & 35 & 80 & 140 & 180 & 200 \\
ECD & 200 & 190 & 165 & 120 & 60 & 20 \\
{\(c_i\)} & 7.5 & 12.5 & 17.5 & 25 & 35 & 50 \\
\end{longtable}

La classe modale est {\([15,20[\)} car sa densité est maximale.

Quartiles et médiane : {\[Q_1=15+5\frac{50-35}{45}\approx16.67,\]} {\[Me=20+10\frac{100-80}{60}\approx23.33,\]} {\[Q_3=30+10\frac{150-140}{40}=32.50.\]} Moyenne : {\[\sum n_ic_i=5075,
\qquad
\bar x=\frac{5075}{200}=25.375.\]} Deuxième moment : {\[\sum n_ic_i^2=154750.\]} Variance : {\[V=\frac{154750}{200}-25.375^2=129.859375.\]} {\[\sigma=\sqrt{129.859375}\approx11.396.\]} {\[CV=\frac{11.396}{25.375}\approx44.91\%.\]} Si 65\% doivent pouvoir payer au moins {\(p\)}, le rang inférieur est {\(35\%\)} : {\[0.35\times200=70.\]} Le rang 70 est dans {\([15,20[\)} : {\[p=15+5\frac{70-35}{45}\approx18.89\text{
euros}.\]} « Satisfait / non satisfait » est une variable qualitative nominale binaire ; un diagramme en barres est adapté.

Tableau complet, avec toutes les fréquences cumulées et les colonnes utilisées pour la moyenne et la variance :

\begin{longtable}[]{@{}llllllllllll@{}}
\toprule\noalign{}
Classe & nᵢ & aᵢ & dᵢ & fᵢ (\%) & ECC & ECD & FCC (\%) & FCD (\%) & cᵢ & nᵢcᵢ & nᵢcᵢ² \\
\midrule\noalign{}
\endhead
\bottomrule\noalign{}
\endlastfoot
{[}5,10{[} & 10 & 5 & 2 & 5 & 10 & 200 & 5 & 100 & 7.5 & 75 & 562.5 \\
{[}10,15{[} & 25 & 5 & 5 & 12.5 & 35 & 190 & 17.5 & 95 & 12.5 & 312.5 & 3906.25 \\
{[}15,20{[} & 45 & 5 & 9 & 22.5 & 80 & 165 & 40 & 82.5 & 17.5 & 787.5 & 13781.25 \\
{[}20,30{[} & 60 & 10 & 6 & 30 & 140 & 120 & 70 & 60 & 25 & 1500 & 37500 \\
{[}30,40{[} & 40 & 10 & 4 & 20 & 180 & 60 & 90 & 30 & 35 & 1400 & 49000 \\
{[}40,60{[} & 20 & 20 & 1 & 10 & 200 & 20 & 100 & 10 & 50 & 1000 & 50000 \\
TOTAL & 200 & --- & --- & 100 & 200 & --- & 100 & --- & --- & 5075 & 154750 \\
\end{longtable}

Les hauteurs de l\textquotesingle histogramme sont les densités : 2 ; 5 ; 9 ; 6 ; 4 ; 1.

\includegraphics[width=0.85\textwidth,height=\textheight]{assets/chapter-7-solutions-fr-chart-01.pdf}

\emph{Cas A --- histogramme des densités / Density histogram}

\includegraphics[width=0.85\textwidth,height=\textheight]{assets/chapter-7-solutions-fr-chart-02.pdf}

\emph{Cas A --- effectifs cumulés / Cumulative counts}

\[\bar x=5075/200=25.375,\quad V=154750/200-(25.375)^2=129.859375.\]

\[\sigma=\sqrt{129.859375}=11.395586,\quad CV=100(11.395586/25.375)=44.9087\%.\]

\[\begin{gathered}
Me=20+\frac{100-80}{60}(10)=23.333333,\\Q_1=15+\frac{50-35}{45}(5)=16.666667,\\Q_3=30+\frac{150-140}{40}(10)=32.5.
\end{gathered}\]

\[P_{35}=15+\frac{70-35}{45}(5)=18.888889\ \mathrm{euros}.\]

Le prix accessible à environ 65 \% des étudiants est le 35e percentile, pas le 65e. Les résultats sont approximatifs car les données sont groupées.

\emph{Préparation CC / partiel --- 110 min}

\hypertarget{cas-b-examen-blanc-inuxe9dit-110-minutes}{%
\subsection{Cas B --- Examen blanc inédit --- 110 minutes}\label{cas-b-examen-blanc-inuxe9dit-110-minutes}}

Cas autonome : les données ci-dessous ne sont pas celles des exercices des chapitres 4 à 6. Présenter toutes les étapes et justifier les arrondis.

\hypertarget{partie-1-double-ruxe9gression-40-min}{%
\subsubsection{Partie 1 --- Double régression (40 min)}\label{partie-1-double-ruxe9gression-40-min}}

\begin{longtable}[]{@{}lllllll@{}}
\toprule\noalign{}
t & 1 & 2 & 3 & 4 & 5 & 6 \\
\midrule\noalign{}
\endhead
\bottomrule\noalign{}
\endlastfoot
x & 5 & 6 & 8 & 9 & 11 & 13 \\
y & 12 & 15 & 18 & 20 & 24 & 29 \\
\end{longtable}

1. Construire le tableau t, x, y, t², x², y², tx, xy ; calculer les deux corrélations et les droites de x sur t puis y sur x. 2. Prévoir x et y pour t=8 ; qualifier la prévision.

\hypertarget{partie-2-suxe9rie-trimestrielle-40-min}{%
\subsubsection{Partie 2 --- Série trimestrielle (40 min)}\label{partie-2-suxe9rie-trimestrielle-40-min}}

\begin{longtable}[]{@{}lllll@{}}
\toprule\noalign{}
Année / Year & Q1 & Q2 & Q3 & Q4 \\
\midrule\noalign{}
\endhead
\bottomrule\noalign{}
\endlastfoot
A & 72 & 88 & 80 & 110 \\
B & 78 & 96 & 88 & 121 \\
\end{longtable}

1. Tracer la série ; calculer les MA4 et les CMA4. 2. Avec la tendance fournie Tₜ=2,8t+78, calculer les huit ratios Yₜ/Tₜ. 3. Calculer les moyennes par trimestre, les normaliser (somme 4) et prévoir t=9 à 12.

\hypertarget{partie-3-indices-30-min}{%
\subsubsection{Partie 3 --- Indices (30 min)}\label{partie-3-indices-30-min}}

\begin{longtable}[]{@{}lllll@{}}
\toprule\noalign{}
Produit / Product & p₀ & q₀ & p₁ & q₁ \\
\midrule\noalign{}
\endhead
\bottomrule\noalign{}
\endlastfoot
A & 10 & 40 & 12 & 45 \\
B & 8 & 60 & 9 & 55 \\
C & 20 & 20 & 22 & 25 \\
\end{longtable}

Construire les quatre agrégats ; calculer les indices de prix et de quantité de Laspeyres, Paasche et Fisher, puis l\textquotesingle indice de valeur ; vérifier la propriété de Fisher.

\hypertarget{corriguxe9-duxe9tailluxe9-1}{%
\subsubsection{Corrigé détaillé}\label{corriguxe9-duxe9tailluxe9-1}}

\hypertarget{partie-1-ruxe9gression}{%
\subsubsection{Partie 1 --- Régression}\label{partie-1-ruxe9gression}}

\begin{longtable}[]{@{}llllllll@{}}
\toprule\noalign{}
t & x & y & t² & x² & y² & t×x & x×y \\
\midrule\noalign{}
\endhead
\bottomrule\noalign{}
\endlastfoot
1 & 5 & 12 & 1 & 25 & 144 & 5 & 60 \\
2 & 6 & 15 & 4 & 36 & 225 & 12 & 90 \\
3 & 8 & 18 & 9 & 64 & 324 & 24 & 144 \\
4 & 9 & 20 & 16 & 81 & 400 & 36 & 180 \\
5 & 11 & 24 & 25 & 121 & 576 & 55 & 264 \\
6 & 13 & 29 & 36 & 169 & 841 & 78 & 377 \\
Σt=21 & 52 & 118 & 91 & 496 & 2510 & 210 & 1115 \\
\end{longtable}

\[\bar t=21/6=3.5,\quad\bar x=52/6=8.666667,\quad\bar y=118/6=19.666667.\]

\[\begin{gathered}
V(t)=91/6-3.5^2=2.916667,\\V(x)=496/6-(52/6)^2=7.555556,\\V(y)=2510/6-(118/6)^2=31.555556.
\end{gathered}\]

\[\begin{gathered}
\operatorname{Cov}(t,x)=210/6-3.5(52/6)=4.666667,\\r_{tx}=\frac{4.666667}{\sqrt{2.916667(7.555556)}}=.994100.
\end{gathered}\]

\[a_1=4.666667/2.916667=1.6,\quad b_1=52/6-1.6(3.5)=3.066667,\quad\widehat x_t=1.6t+3.066667.\]

\[\begin{gathered}
\operatorname{Cov}(x,y)=1115/6-(52/6)(118/6)=15.388889,\\r_{xy}=\frac{15.388889}{\sqrt{7.555556(31.555556)}}=.996635.
\end{gathered}\]

\[a_2=15.388889/7.555556=2.036765,\quad b_2=118/6-a_2(52/6)=2.014706.\]

\[\begin{gathered}
\widehat x_8=1.6(8)+3.066667=15.866667,\\\widehat y_8=2.036765(15.866667)+2.014706=34.331373.
\end{gathered}\]

La période 8 dépasse les observations 1 à 6 : c\textquotesingle est une extrapolation. Corrélation n\textquotesingle implique pas causalité.

\hypertarget{partie-2-suxe9rie-temporelle}{%
\subsubsection{Partie 2 --- Série temporelle}\label{partie-2-suxe9rie-temporelle}}

\[MA_{2.5}=(72+88+80+110)/4=87.5,\quad MA_{3.5}=(88+80+110+78)/4=89.\]

\begin{longtable}[]{@{}lll@{}}
\toprule\noalign{}
Fenêtre / Window & MA4 & Période médiane / Midperiod \\
\midrule\noalign{}
\endhead
\bottomrule\noalign{}
\endlastfoot
1--4 & 87.5 & 2.5 \\
2--5 & 89 & 3.5 \\
3--6 & 91 & 4.5 \\
4--7 & 93 & 5.5 \\
5--8 & 95.75 & 6.5 \\
\end{longtable}

\[CMA_3=(87.5+89)/2=88.25,\quad CMA_4=(89+91)/2=90.\]

\begin{longtable}[]{@{}p{0.32\linewidth}p{0.62\linewidth}@{}}
\toprule\noalign{}
t & CMA4 \\
\midrule\noalign{}
\endhead
\bottomrule\noalign{}
\endlastfoot
3 & 88.25 \\
4 & 90 \\
5 & 92 \\
6 & 94.375 \\
\end{longtable}

\begin{longtable}[]{@{}llll@{}}
\toprule\noalign{}
t & Yₜ & Tₜ=2.8t+78 & Yₜ/Tₜ \\
\midrule\noalign{}
\endhead
\bottomrule\noalign{}
\endlastfoot
1 & 72 & 80.8 & 0.891089 \\
2 & 88 & 83.6 & 1.052632 \\
3 & 80 & 86.4 & 0.925926 \\
4 & 110 & 89.2 & 1.233184 \\
5 & 78 & 92 & 0.847826 \\
6 & 96 & 94.8 & 1.012658 \\
7 & 88 & 97.6 & 0.901639 \\
8 & 121 & 100.4 & 1.205179 \\
\end{longtable}

\[\bar S_q=(S_q+S_{q+4})/2,\qquad S'_q=\frac{4\bar S_q}{\sum_{k=1}^4\bar S_k}.\]

\begin{longtable}[]{@{}lll@{}}
\toprule\noalign{}
Quarter & Moyenne / Mean & S\textquotesingle{} normalisé / Normalised \\
\midrule\noalign{}
\endhead
\bottomrule\noalign{}
\endlastfoot
Q1 & 0.869458 & 0.861902 \\
Q2 & 1.032645 & 1.023671 \\
Q3 & 0.913783 & 0.905841 \\
Q4 & 1.219182 & 1.208586 \\
TOTAL & 4.035067 & 4 \\
\end{longtable}

\begin{longtable}[]{@{}llll@{}}
\toprule\noalign{}
t & Tₜ & S\textquotesingle{} & Y prévu / Forecast \\
\midrule\noalign{}
\endhead
\bottomrule\noalign{}
\endlastfoot
9 & 103.2 & 0.861902 & 88.9482 \\
10 & 106 & 1.023671 & 108.5091 \\
11 & 108.8 & 0.905841 & 98.5555 \\
12 & 111.6 & 1.208586 & 134.8782 \\
\end{longtable}

\includegraphics[width=0.85\textwidth,height=\textheight]{assets/chapter-7-solutions-fr-chart-03.pdf}

\emph{Cas B --- série trimestrielle / Quarterly series}

Prévisions t=9 à 12 : 88,95 ; 108,51 ; 98,56 ; 134,88.

\hypertarget{partie-3-indices}{%
\subsubsection{Partie 3 --- Indices}\label{partie-3-indices}}

\begin{longtable}[]{@{}lllll@{}}
\toprule\noalign{}
Produit / Product & p₀q₀ & p₁q₀ & p₀q₁ & p₁q₁ \\
\midrule\noalign{}
\endhead
\bottomrule\noalign{}
\endlastfoot
A & 400 & 480 & 450 & 540 \\
B & 480 & 540 & 440 & 495 \\
C & 400 & 440 & 500 & 550 \\
Σ & 1280 & 1460 & 1390 & 1585 \\
\end{longtable}

\[L_P=100(1460/1280)=114.0625,\quad P_P=100(1585/1390)=114.028777.\]

\[L_Q=100(1390/1280)=108.59375,\quad P_Q=100(1585/1460)=108.561644.\]

\[F_P=\sqrt{L_PP_P}=114.045637,\quad F_Q=\sqrt{L_QP_Q}=108.577696.\]

\[I_V=100(1585/1280)=123.828125,\quad F_PF_Q/100=123.828125.\]

La valeur totale augmente de 23,83 \%, sous l\textquotesingle effet conjoint des prix et des quantités.

\end{document}
